\subsection{Comparison with the constant graphon}
\label{sec:constant-graphon-comparison}

The constant graphon $W\equiv r_h$ has $H$-density $r_h^m$ and is therefore
feasible for the same constraint as $W_h$, where $W_h$ is the member of the
analytic rank-one KKT family constructed in
\Cref{lem:rank-one-kkt-family}.  We will prove that
\[
  I_{p_h}(W_h)<I_{p_h}(W\equiv r_h)=J_{p_h}(r_h),
\]
and compute the asymptotic expansion of the
relative entropy difference for all sufficiently small $h>0$.  At $h=0$, the
two graphons coincide, so their $I_{p_h}$-values are equal.  The symmetries in
\Cref{lem:rank-one-kkt-family} make the relative entropy difference
an even analytic function of $h$, and the calculation below shows that its
leading term is of order $h^4$.  The sign of the difference for small $h>0$ is therefore
determined by the coefficient of $h^4$.  To compute that coefficient, we
first record the expansions of the parameters that will be used below.
Write
\[
  \ell(z):=\log\frac{1-z}{z}, \qquad
  \ell_*:=\ell(p_*)=\frac d{d-1}-\log(d-1), \qquad
  \Lambda_h:=\ell(p_h)-\ell_*.
\]
Here, $\Lambda_h$ represents the slope of the affine difference
\begin{equation}\label{eq:entropy-parameter-shift}
  J_{p_h}(z)-J_{p_*}(z)=\Lambda_h z+J_{p_h}(0)-J_{p_*}(0).
\end{equation}

\begin{lemma}[Expansions along the rank-one KKT family]
\label{lem:rank-one-parameter-expansions}
The analytic family constructed in
\Cref{lem:rank-one-kkt-family} satisfies
\begin{equation}\label{eq:rank-one-parameter-expansions}
\begin{aligned}
  u_h&=u_*+\frac{5-3d}{6u_*}h^2+O_d(h^4),
  &&&
  \Lambda_h&=\frac{2d^3}{3(d-1)}h^2+O_d(h^4),\\
  \alpha_h&=\frac12+\frac{2d^2u_*}{3(d-1)}h+O_d(h^3),
  &&&
  r_h&=r_*-\frac{2(4d-1)}3h^2+O_d(h^4).
\end{aligned}
\end{equation}
\end{lemma}

These expansions immediately locate the family relative to
$P_*=(p_*,r_*)$.  For all sufficiently small $h>0$, the expansion of $r_h$
gives $r_h<r_*$.  Moreover, since $\ell$ is strictly decreasing,
$\Lambda_h>0$ implies $p_h<p_*$.

\begin{proof}[Proof outline]
The symmetries in \Cref{lem:rank-one-kkt-family} show that
$u_h,\gamma_h$, and $\Lambda_h$ have expansions in even powers of $h$, while
$\alpha_h-1/2$ has an expansion in odd powers.  Let $U_2, G_2, L_2$
denote the quadratic coefficients of $u_h,\gamma_h$, and $\Lambda_h$,
respectively, and substitute the expansions of $s_h^2,s_ht_h,t_h^2$ into
\eqref{eq:three-value-kkt}.
The middle equation at order $h^2$ relates $G_2$ and $L_2$,
and either outer equation at order $h^3$ determines $G_2$,
and the comparison at order $h^4$ relates $U_2$ and $G_2$.
Solving these three equations gives the first two expansions in
\eqref{eq:rank-one-parameter-expansions}.

Now we determine the expansions of $\alpha_h$ and $r_h$.
Since $R(h,z)$ is even in $h$, its Taylor expansion gives
\[
  R(h,z)=k_3+k_4(z-r_*)+O_d(h^2)
  \quad (z\in[s_h^2,t_h^2]),
  \qquad
  k_4:=\frac{5d^6(d-2)}{24(d-1)^3}.
\]
Using this expansion and integrating the factorization
\eqref{eq:mh-derivative-factorization} over $[s_h^2,s_ht_h]$ and
$[s_ht_h,t_h^2]$ gives the two differences in \eqref{eq:block-proportion-formula}
through order $h^5$.  Their asymmetry at order $h^5$ determines the linear
term of $\alpha_h$.  The expansions of $q_h$ and $r_h$ then follow from
their definitions.
Full details are given in \Cref{app:rank-one-parameter-expansions}.
\end{proof}

We now turn to the comparison stated at the beginning of this subsection.

\begin{lemma}[Comparison with the constant graphon]
\label{lem:constant-graphon-comparison}
For sufficiently small $h>0$,
\[
  e(W_h)=r_h-(d-1)h^2+O_d(h^4),
  \qquad
  I_{p_h}(W_h)=J_{p_h}(r_h)-\frac{d^3}{3}h^4+O_d(h^6).
\]
In particular, the constant graphon
$W\equiv r_h$ is feasible but has strictly larger $I_{p_h}$-value,
and consequently $p_h<\pc(r_h)$ for all sufficiently small $h>0$.
\end{lemma}

The affine identity \eqref{eq:entropy-parameter-shift} separates the change in $p$
from the comparison at $p_*$:
\[
  I_{p_h}(W_h)-J_{p_h}(r_h)
  =\{I_{p_*}(W_h)-J_{p_*}(r_h)\}+\Lambda_h\{e(W_h)-r_h\}.
\]
The second term is evaluated using \Cref{lem:rank-one-parameter-expansions}.
For the first difference, we return to the supporting-line geometry of
\Cref{thm:scalar-lz-boundary}.  Its proof obtains the quadratic separation
in \eqref{eq:contact-quadratic-separation} by Taylor-expanding the difference between
$\varphi_{\pc(r),d}(x)=J_{\pc(r)}(x^{1/d})$ and its tangent line $\ell_r(x)$
near the two contact points.  We use the same approach at the exceptional
endpoint, where these contacts coincide.

By the continuous extensions in \Cref{thm:scalar-lz-boundary},
$\pc(r)\to p_*$ and $\sm(r)\to r_*$ as $r\to r_*$.  The tangent lines
$\ell_r$ therefore extend continuously to the tangent of
$\varphi_{p_*,d}$ at $x=r_*^d$:
\[
  \ell_{r_*}(x)=J_{p_*}(r_*)+\beta_d(x-r_*^d).
\]
Here $\beta_d$ denotes the slope of this tangent, which the chain rule gives as
\[
  \beta_d:=\varphi_{p_*,d}'(r_*^d)
  =\frac{J_{p_*}'(r_*)}{d r_*^{d-1}}
  =\frac{\gamma_*}{d}
  =\frac{d^{d-1}}{(d-1)^d}.
\]
Following \eqref{eq:supporting-cost-gap}, define the \emph{supporting cost gap} by
\begin{equation}\label{eq:endpoint-supporting-gap}
  \Gamma_d(z)
  :=\varphi_{p_*,d}(z^d)-\ell_{r_*}(z^d)
  =J_{p_*}(z)-J_{p_*}(r_*)-\beta_d\{z^d-r_*^d\}
  \qquad(0\le z\le1).
\end{equation}
Thus, $\Gamma_d$ is the continuous extension to $r=r_*$ of the function
$g_r$ used there.  Its Taylor expansion describes the separation from
the tangent at their common contact point; as shown below, the leading term is quartic.

By \eqref{eq:endpoint-supporting-gap}, we can express the cost difference
at $p_*$ in terms of $\Gamma_d$ as
\[
  I_{p_*}(W_h)-J_{p_*}(r_h)
  =\int\Gamma_d(W_h)-\Gamma_d(r_h)+\beta_d\left(\int W_h^d-r_h^d\right).
\]
Adding the affine correction from $p_*$ to $p_h$ gives
\begin{equation}\label{eq:constant-comparison-decomposition}
  I_{p_h}(W_h)-J_{p_h}(r_h)
  =\int\Gamma_d(W_h)-\Gamma_d(r_h)
  +\beta_d\left(\int W_h^d-r_h^d\right)
  +\Lambda_h\{e(W_h)-r_h\}.
\end{equation}
Since $\int W_h^d=(\int f_h^d)^2=q_h^2=r_h^d$, the $d$-th moment term vanishes, and
\begin{equation}\label{eq:constant-comparison-identity}
  I_{p_h}(W_h)-J_{p_h}(r_h)
  =\int\Gamma_d(W_h)-\Gamma_d(r_h) +\Lambda_h\{e(W_h)-r_h\}.
\end{equation}
The proof first evaluates $\Gamma_d$ by its quartic expansion,
then expands all the parameters using \Cref{lem:rank-one-parameter-expansions}.

\begin{proof}
Tangency gives $\Gamma_d(r_*)=\Gamma_d'(r_*)=0$.  At the exceptional
endpoint, direct differentiation further gives
$\Gamma_d''(r_*)=\Gamma_d'''(r_*)=0$ and
$\Gamma_d^{(4)}(r_*)=d^5/(d-1)^2>0$.
Taylor's theorem therefore yields, as $z\to r_*$,
\begin{equation}\label{eq:endpoint-gap-expansion}
  \Gamma_d(z)
  =\frac{d^5}{24(d-1)^2}(z-r_*)^4+O_d(|z-r_*|^5).
\end{equation}
By the parameter expansions in \eqref{eq:rank-one-parameter-expansions},
\[
  s_h^2-r_*=-2u_*h+O_d(h^2),
  \qquad
  t_h^2-r_*=2u_*h+O_d(h^2),
  \qquad
  s_ht_h-r_*=O_d(h^2),
  \qquad
  r_h-r_*=O_d(h^2).
\]
Each of the four rectangles determined by $A_h$ and $A_h^c$ in
$[0,1]^2$ has area $1/4+O_d(h)$.
Thus the quartic expansion and $u_*^2=(d-1)/d$ yield
\begin{equation}\label{eq:family-gap-expansion}
  \int\Gamma_d(W_h)=\frac{d^3}{3}h^4+O_d(h^6),
  \qquad
  \Gamma_d(r_h)=O_d(h^8).
\end{equation}
The first remainder improves from $O_d(h^5)$ to $O_d(h^6)$
because $\int\Gamma_d(W_h)$ is an even analytic function of $h$.
Finally, $e(W_h)=\{u_h+h(1-2\alpha_h)\}^2$ and the parameter expansions give
\[
  e(W_h)-r_h=-(d-1)h^2+O_d(h^4),
  \qquad
  \Lambda_h\{e(W_h)-r_h\}=-\frac{2d^3}{3}h^4+O_d(h^6).
\]
Substituting these expansions into \eqref{eq:constant-comparison-identity}
proves the lemma.
\end{proof}


